\section{Scope and interpretation}
\paragraph{A full-support prior with an explicit inductive preference.}
Every finite parameter vector has positive prior density. Large internal gains reduce endpoint variances instead of being forbidden. Theorem~\ref{thm:prior} describes the resulting probabilities of function neighborhoods. Theorem~\ref{thm:energy-characterization} identifies the limiting minimum normalized energy with twice minimal path variation. The conditional Gaussian normalizers, including the tilted internal marginal, are included in the probability proof.

\paragraph{Why path variation is a natural quantity.}
Positive neuron rescalings leave a ReLU function and its path products unchanged. Path regularizers and this rescaling symmetry are established tools in neural-network optimization~\cite{path-sgd}. Total path variation also controls width-independent statistical complexity~\cite{barron-klusowski}. The mixed version here uses Euclidean first-row norms because the input domain is a Euclidean unit ball. Our contribution is the explicit Gaussian construction whose function probabilities and ordinary posterior satisfy the stated bounds; we do not claim that the path complexity measure itself is new.

\paragraph{What changed from independent Gaussian weights.}
The base internal coordinates have variance $N^{-\gamma}$, and the base endpoint variances are $N/\omega$ and $1/(\omega N)$. The soft gain penalty changes the internal marginal and gives gain-normalized conditional endpoint variances. At $\gamma=1$, $\omega$ is exactly the Gaussian parameter count. No additional width threshold or logarithmic precision correction is needed. The theorem does not prove the same path-variation bound for the unchanged independent-coordinate prior.

\paragraph{Other internal Gaussian scales.}
The proof also allows arbitrary positive internal precisions $\alpha_\ell$. Use base internal variance $1/\alpha_\ell$ and set
$\omega=N(d+1)+N\sum_{\ell=2}^H\max\{N,\alpha_\ell\}$.
The gain moment still has $D_{\mathrm{int}}/\omega\le1$, the total dimension ratio is at most one, and the cloned internal mean cost is at most $C/N$. Every other argument is unchanged. This is a calibrated family of softly regularized priors, not an assertion about arbitrary unchanged independent Gaussian exponent tuples.

\paragraph{Small width versus small path variation.}
Every finite ReLU function has a flow-balanced representation at its existing width and depth. Nevertheless a small network can have large path variation, and the theorem does not give unrestricted polynomial dependence on width and depth. The quadratic-width, linear-depth corollary requires bounded endpoint norms and bounded product of internal absolute spectral gains. No such teacher bound is used to choose the prior or likelihood.

\paragraph{Noise, computation, and verification.}
The labels may follow any bounded IID law, without a network target or Gaussian noise assumption. The data-dependent likelihood yields faster rates at smaller comparison risk and predicts by averaging raw network outputs. The theorem provides neither efficient evaluation of free energy nor a gradient-descent guarantee. These conventional proofs have not been formalized in Lean.
